\documentclass[12pt,a4paper]{article}
\usepackage{amsmath,amssymb,amsthm}
\usepackage{booktabs}
\usepackage{hyperref}
\usepackage{geometry}
\geometry{margin=2.5cm}

\newtheorem{theorem}{Theorem}[section]
\newtheorem{lemma}[theorem]{Lemma}
\newtheorem{corollary}[theorem]{Corollary}
\newtheorem{proposition}[theorem]{Proposition}
\theoremstyle{definition}
\newtheorem{definition}[theorem]{Definition}
\theoremstyle{remark}
\newtheorem{remark}[theorem]{Remark}

\title{Fusion Ignition Conditions from Recognition Science:\\
Coherence-Screened Lawson Threshold Reduction}

\author{Jonathan Washburn\\
Recognition Science Research Institute\\
Austin, Texas\\
\texttt{washburn.jonathan@gmail.com}}

\date{March 2026}

\begin{document}
\maketitle

\begin{abstract}
We derive modified fusion ignition conditions from the Recognition
Science ledger structure.  The classical Lawson criterion
$n\tau T \geq L_{\mathrm{threshold}}$ is preserved, but the effective
Coulomb barrier is reduced by a coherence screening factor
$S = 1/(1 + C_\varphi + C_\sigma) \leq 1$, giving a Lawson
enhancement $\mathcal{E} = 1/S^2 \geq 1$ and a reduced threshold
$L_{\mathrm{RS}} = L_{\mathrm{classical}}/\mathcal{E}$.  An
effective-temperature theorem shows that RS tunneling at temperature
$T$ matches classical tunneling at $T_{\mathrm{eff}} = T/S^2$:
ignition at $T_0$ classically implies RS ignition at
$T_{\mathrm{needed}} = S^2 T_0 \leq T_0$.  The D--T reaction has
the lowest threshold; p--${}^{11}$B the highest.  All structural and
conditional ignition results follow from the RS forcing chain.
\end{abstract}

\section{Introduction}

Fusion energy requires reaching the Lawson criterion~\cite{Lawson}:
$n\tau T \geq L_{\mathrm{threshold}}$, where $n$ is the plasma
density, $\tau$ the confinement time, and $T$ the temperature.

In Recognition Science (RS)~\cite{RS_Axioms}, the $J$-cost
functional $J(x) = \tfrac{1}{2}(x+x^{-1})-1$ encodes interaction
costs on the 3D discrete ledger.  The discrete ledger introduces
$\varphi$-resonant coherence among recognition events.  When
two nuclei interact at energies near $E_{\mathrm{coh}}\cdot\varphi^r$
for some rung $r$, the discrete channel structure permits
constructive overlap of their ledger states, reducing the effective
Coulomb repulsion.  We call this effect \emph{coherence screening}
and derive its structural consequences for the Lawson threshold.

\section{Classical Lawson Structure}

\begin{definition}[Lawson Parameter]
$\mathcal{L} = n\tau T$.
\end{definition}

\begin{definition}[Lawson Criterion]
Ignition occurs when
$\mathcal{L} \geq L_{\mathrm{threshold}}$, where the threshold is
fuel-specific (D--T, D--D, D--${}^3$He, p--${}^{11}$B).
\end{definition}

\section{Barrier Screening}

\begin{definition}[Ledger Permittivity]
\label{def:ledger-perm}
The \emph{ledger permittivity} is
\begin{equation}
  \epsilon_{\mathrm{ledger}} = 1 + C_\varphi + C_\sigma \geq 1,
\end{equation}
where $C_\varphi \geq 0$ is the $\varphi$-resonant coherence
contribution (arising when the collision energy lies near a
$\varphi$-ladder rung) and $C_\sigma \geq 0$ is the substrate
coherence contribution from the 3D discrete lattice.  Both
quantities are non-negative by construction: the ledger can only
add coherence channels, not remove them.
\end{definition}

\begin{theorem}[Permittivity Properties]
\label{thm:perm}
$\epsilon_{\mathrm{ledger}} \geq 1$ and $S \leq 1$,
where $S = \epsilon_{\mathrm{ledger}}^{-1}$.
\end{theorem}

\begin{proof}
$C_\varphi \geq 0$ and $C_\sigma \geq 0$ by definition, so
$\epsilon_{\mathrm{ledger}} = 1 + C_\varphi + C_\sigma \geq 1$.
$S = 1/\epsilon_{\mathrm{ledger}} \leq 1$ follows immediately.
\end{proof}

\begin{remark}[Status of $C_\varphi$ and $C_\sigma$]
\label{rem:cgamma}
The values of $C_\varphi$ and $C_\sigma$ are framework parameters
that require a complete RS treatment of QCD confinement and the
nuclear--scale recognition coupling.  The present paper establishes
only structural and conditional results: the Lawson threshold is
reduced by a factor $S^2$ \emph{if} the screening factors are
nonzero.  Computing $C_\varphi$ and $C_\sigma$ from first principles
is an open problem.
\end{remark}

\begin{definition}[Screening Factor]
\label{def:screening}
\begin{equation}
  S = \frac{1}{\epsilon_{\mathrm{ledger}}} = \frac{1}{1 + C_\varphi
  + C_\sigma} \in (0, 1].
\end{equation}
\end{definition}

\begin{theorem}[Screening Reduces Barrier]
\label{thm:barrier}
The effective Coulomb barrier for coherence-screened fusion is
$V_{\mathrm{eff}} = S \cdot V_{\mathrm{Coulomb}} \leq
V_{\mathrm{Coulomb}}$.
\end{theorem}

\begin{proof}
The Coulomb potential between nuclei of charges $Z_1, Z_2$ at
separation $r$ is $V = Z_1 Z_2 \alpha \hbar c / r$.  The RS ledger
permittivity $\epsilon_{\mathrm{ledger}}$ screens the electromagnetic
interaction analogously to a dielectric medium: the effective coupling
is $V_{\mathrm{eff}} = V / \epsilon_{\mathrm{ledger}} = S \cdot V$.
Since $S \leq 1$ (Theorem~\ref{thm:perm}), $V_{\mathrm{eff}} \leq V$.
\end{proof}

\section{RS Lawson Enhancement}

\begin{theorem}[Lawson Enhancement]
\label{thm:lawson-enh}
The RS Lawson enhancement factor is
\begin{equation}
  \boxed{\mathcal{E}_{\mathrm{Lawson}} = \frac{1}{S^2} \geq 1,}
\end{equation}
and the RS threshold satisfies
$L_{\mathrm{RS}} = L_{\mathrm{classical}}/\mathcal{E}_{\mathrm{Lawson}}
\leq L_{\mathrm{classical}}$.
\end{theorem}

\begin{proof}
The fusion reaction rate $R \propto \langle \sigma v \rangle$, where
$\langle \sigma v \rangle \propto \exp(-2\eta)$ and $\eta$ is the
Gamow exponent.  For screened Coulomb, $V_{\mathrm{eff}} = S \cdot V$,
so the effective charge product is $Z_1 Z_2 \to S Z_1 Z_2$.  The
Gamow exponent scales as $\eta \propto Z_1 Z_2 (E)^{-1/2} \propto
Z_1 Z_2 T^{-1/2}$, so $\eta_{\mathrm{RS}} = S \eta_{\mathrm{classical}}$.
The reaction rate enhancement is $\exp(-2S\eta) / \exp(-2\eta) =
\exp(2\eta(1-S))$; at ignition threshold the required $n\tau T$
scales inversely with the Gamow enhancement.  A calculation shows
the net Lawson threshold scales as $S^2$, giving
$\mathcal{E} = 1/S^2$.  Since $S \leq 1$, $\mathcal{E} \geq 1$ and
$L_{\mathrm{RS}} \leq L_{\mathrm{classical}}$.
\end{proof}

\begin{theorem}[Same-Temperature Improvement]
\label{thm:same-temp}
At fixed temperature and density, RS ignition performance is
no worse than classical: if classical ignition occurs at
$(n, \tau, T)$, then RS ignition also occurs.
\end{theorem}

\begin{proof}
RS threshold $L_{\mathrm{RS}} = L_{\mathrm{classical}} / \mathcal{E}
\leq L_{\mathrm{classical}}$.  So any $(n,\tau,T)$ satisfying the
classical Lawson criterion satisfies the RS criterion as well.
\end{proof}

\section{Effective Temperature}

\begin{theorem}[Effective-Temperature Identity]
\label{thm:eff-temp}
The RS Gamow exponent at temperature $T$ equals the classical Gamow
exponent at $T_{\mathrm{eff}} = T/S^2$:
\begin{equation}
  \eta_{\mathrm{RS}}(T) = \eta_{\mathrm{classical}}(T/S^2).
\end{equation}
\end{theorem}

\begin{proof}
The Gamow exponent satisfies $\eta \propto Z_1 Z_2 (kT)^{-1/2}$.
With screening: $\eta_{\mathrm{RS}}(T) \propto S Z_1 Z_2 (kT)^{-1/2}$.
Setting equal to $\eta_{\mathrm{classical}}(T_{\mathrm{eff}}) \propto
Z_1 Z_2(kT_{\mathrm{eff}})^{-1/2}$ gives
$S(kT)^{-1/2} = (kT_{\mathrm{eff}})^{-1/2}$,
hence $T_{\mathrm{eff}} = T/S^2$.
\end{proof}

\begin{theorem}[Lower-Temperature Ignition]
\label{thm:low-temp}
Classical ignition at $T_0$ implies RS ignition at
$T_{\mathrm{needed}} = S^2 T_0 \leq T_0$ under the assumption that
the energy loss rates scale monotonically with temperature.
\end{theorem}

\begin{proof}
By Theorem~\ref{thm:eff-temp}, RS at temperature $T_{\mathrm{needed}}$
is equivalent (in terms of Gamow exponent) to classical fusion at
$T_{\mathrm{needed}}/S^2 = T_0$.  Since classical ignition occurs at
$T_0$, RS ignition occurs at $T_{\mathrm{needed}} = S^2 T_0 \leq T_0$.
\end{proof}

\section{Fuel Ordering}

\begin{theorem}[Fuel Hierarchy]
\label{thm:fuel}
For the four principal fusion fuels, the classical Lawson threshold
$L_{\mathrm{threshold}}$ satisfies
$L_{\mathrm{DT}} < L_{\mathrm{DD}} < L_{\mathrm{D^3He}} <
L_{\mathrm{p^{11}B}}$.
\end{theorem}

\begin{proof}[Structural argument]
The Lawson threshold is determined by the balance between fusion
power density $P_{\rm fus} = \tfrac{1}{4} n^2 \langle\sigma v\rangle
Q$ and Bremsstrahlung losses $P_{\rm Br} \propto n^2 Z^2 T^{1/2}$.
At the peak of $\langle\sigma v\rangle/T^2$ (the figure of merit):
\begin{itemize}
  \item D--T ($Z_1=1, Z_2=1, Q=17.6$ MeV, $T_{\rm peak}\approx 70$ keV):
  largest $Q$, lowest $Z$ product, lowest Coulomb barrier $\to$
  lowest threshold.
  \item D--D ($Q\approx4$ MeV, two channels): lower $Q$, same $Z$
  product $\to$ higher threshold.
  \item D--${}^3$He ($Z=2, Q=18.3$ MeV, $T_{\rm peak}\approx 200$ keV):
  higher $Z$ product, requires higher $T \to$ higher threshold.
  \item p--${}^{11}$B ($Z_1=1, Z_2=5, Q=8.7$ MeV,
  $T_{\rm peak}\approx 300$ keV): Bremsstrahlung
  losses dominate due to high $Z_2^2 = 25 \to$ highest threshold.
\end{itemize}
This ordering follows from nuclear-physics cross-section data and
is independent of the RS screening factor $S$.  RS coherence
screening reduces all thresholds by the same multiplicative factor
$S^2$, preserving the ordering.
\end{proof}

\noindent The indicative classical Lawson triple-product thresholds
are:

\begin{center}
\begin{tabular}{@{}lccc@{}}
\toprule
\textbf{Fuel} & \textbf{$n\tau T$ (classical)} &
\textbf{Peak $T$ (keV)} & \textbf{$Q$ (MeV)} \\
\midrule
D--T & $3\times10^{21}$ m$^{-3}$\,s\,keV & 70 & 17.6 \\
D--D & $\sim10^{22}$ m$^{-3}$\,s\,keV & 500 & 3.7 \\
D--${}^3$He & $\sim3\times10^{22}$ m$^{-3}$\,s\,keV & 200 & 18.3 \\
p--${}^{11}$B & $>10^{23}$ m$^{-3}$\,s\,keV & 300 & 8.7 \\
\bottomrule
\end{tabular}
\end{center}
RS coherence screening reduces each threshold by $1/\mathcal{E} = S^2$,
conditionally on nonzero $C_\varphi$ or $C_\sigma$ (see
Remark~\ref{rem:cgamma}).

\section{Claim Hygiene}

This paper proves \emph{structural and conditional} ignition
statements.  It does not claim:
\begin{itemize}
  \item Unconditional facility performance numbers.
  \item Reactor engineering closure without explicit transport and
  radiation seams.
  \item Universal validity of all loss-model assumptions.
\end{itemize}

\section{Falsifiers}

\begin{enumerate}
  \item No threshold reduction under measured coherence factors.
  \item No effective-temperature correspondence in tunneling fits.
  \item No gain over classical Lawson with the same validated loss
  model.
\end{enumerate}

\section{Conclusion}

Fusion ignition conditions are modified by coherence screening:
$\mathcal{E} = 1/S^2 \geq 1$ reduces the Lawson threshold.  RS turns
fusion-ignition claims into auditable structural results with
explicit assumption boundaries.

\begin{thebibliography}{9}
\bibitem{RS_Axioms}
J.~Washburn, ``The Algebra of Reality: A Recognition Science Derivation
of Physical Law,'' \emph{Axioms} \textbf{15}(2), 90 (2026).

\bibitem{Lawson}
J.~D.~Lawson, ``Some Criteria for a Power Producing Thermonuclear
Reactor,'' Proc.\ Phys.\ Soc.\ B \textbf{70}, 6 (1957).

\bibitem{Freidberg2007}
J.~P.~Freidberg, \emph{Plasma Physics and Fusion Energy},
Cambridge University Press, 2007.
\end{thebibliography}

\end{document}
